\documentclass[spanish,11pt,a4paper]{article}
\usepackage[margin=1in]{geometry}
\usepackage{fontspec}
\setmainfont{lmroman10-regular.otf}[BoldFont=lmroman10-bold.otf,ItalicFont=lmroman10-italic.otf,BoldItalicFont=lmroman10-bolditalic.otf]
\setsansfont{lmsans10-regular.otf}[BoldFont=lmsans10-bold.otf]
\setmonofont{lmmono10-regular.otf}[Scale=MatchLowercase]
\newlength{\LeanCharWidth}
\usepackage{amsmath,amssymb}
\usepackage[spanish,es-noshorthands]{babel}
\usepackage{microtype,parskip}
\usepackage{graphicx,float,caption}
\usepackage{longtable,booktabs,array,calc}
\usepackage{xcolor}
\usepackage{xurl}
\usepackage{hyperref,bookmark}
\usepackage{etoolbox,needspace}
\setcounter{secnumdepth}{-\maxdimen}
\setlength{\emergencystretch}{3em}
\providecommand{\tightlist}{\setlength{\itemsep}{0pt}\setlength{\parskip}{0pt}}
\newcounter{none}
\AtBeginEnvironment{longtable}{\Needspace{8\baselineskip}\small}
\pretocmd{\section}{\Needspace{5\baselineskip}}{}{}
\pretocmd{\subsection}{\Needspace{4\baselineskip}}{}{}
\pretocmd{\subsubsection}{\Needspace{4\baselineskip}}{}{}
\captionsetup{font=small,labelfont=bf}
\urlstyle{same}
\AtBeginDocument{\special{pdf:minorversion 7}}

\hypersetup{pdflang=es-CL,colorlinks=true,urlcolor=blue,linkcolor=blue,pdftitle={Cuaderno de demostraciones de Recuperación de descomposiciones en triángulos con anfitrión tras borrados aleatorios en grafos split},pdfauthor={Juan Pablo Traverso Gianini}}
\title{Cuaderno de demostraciones de Recuperación de descomposiciones en triángulos con anfitrión tras borrados aleatorios en grafos split}
\author{Juan Pablo Traverso Gianini\\\small Investigador independiente, Santiago, Chile\\\small\href{mailto:jtraverso@gmail.com}{jtraverso@gmail.com}\quad\href{https://orcid.org/0009-0003-6068-4096}{ORCID 0009-0003-6068-4096}}
\date{Preprint 1.0 \enspace·\enspace 24 de noviembre de 2026}
\begin{document}\raggedbottom\maketitle
\begin{center}Artículo III de la serie sobre brechas de partición en cliques en grafos split.\end{center}
\noindent DOI \url{https://doi.org/10.5281/zenodo.24494747}\\
Repositorio de la serie \url{https://github.com/jtraverso/split-clique-partition-gap}
\subsection{De los triángulos asignados a la brecha}\label{de-los-triuxe1ngulos-asignados-a-la-brecha}

Sea \(J\) el conjunto de pares requeridos y \(B\) el grafo residual sobre la clique. Al asignar una base se usan dos aristas cruzadas distintas. Una asignación propia usa cada arista cruzada a lo sumo una vez. Si la asignación es completa, los dos costes de partición en (7) tienen el término común \(|E(G)|-2|J|\). Restar el coste fraccionario del entero da \(F(B)-I(B)\). Al borrar aristas cruzadas se modifica el término común y se conserva la independencia en \(B\) de las vecindades de los anfitriones. Una nueva asignación de todos los pares da por tanto la misma diferencia. Este es el paso de transferencia en el Teorema 3.

\subsection{Recuperación de las cuatro hipótesis en un mismo soporte}\label{recuperaciuxf3n-de-las-cuatro-hipuxf3tesis-en-un-mismo-soporte}

En cada lado de un bloque bipartito, el factor debe llevar \(d^2\) unidades. En un corte cuyo lado de la fuente contiene a \(X\) y a los anfitriones fuera de \(Y\), la capacidad disponible es \[d(d-|X|)+e(X,Y)+\sum_{h\notin Y}a_h.\] Al usar \(\sum_h a_h=d^2\) y comparar con \(d^2\) se obtiene (11). Usar los mismos grados de anfitrión en ambos lados permite emparejar las incidencias seleccionadas en cada anfitrión. En el bloque clique, los grados del núcleo son \(d-1\) y los grados de los anfitriones deben ser pares. El apéndice B controla todos los cortes antes de elegir los factores.

Los pesos iniciales de triángulo saturan las aristas cruzadas seleccionadas. Los errores restantes están en las bases y tienen márgenes nulos. Las correcciones de rectángulos de la sección 5 conservan las cargas cruzadas. En el caso de la clique, el coeficiente de \(D_{ij}\) en (12) resulta de sumar cuatro contribuciones \[2(d-2)(d-3)+4+2(d-3)+2(d-3)=2(d-1)(d-2).\] El reparto de cada corrección entre anfitriones comunes mantiene positivos los pesos corregidos. El argumento entero de la sección 5 proporciona H2. Los conteos de vecinos comunes tras el recorte dan H3 e incluyen los pares sin incidencia superviviente.

Para la bisagra que comparte una arista cruzada, los triángulos raíz son \(uhv\) y \(uhw\). Las siete aristas auxiliares son \(ku,kv,kw,kz,vz,wz,hz\). Al añadir \(uv,hv\) se obtienen \(kuv,hvz,kwz\). Al añadir, en cambio, \(uw,hw\) se obtienen \(kuw,hwz,kvz\). Las nueve aristas aparecen una vez en cada partición. Los vértices raíz pertenecen al soporte, y las aristas auxiliares son disjuntas de las de ambos triángulos raíz. Los papeles de \(k\) y \(z\) se recuperan a partir de los únicos vecinos auxiliares de \(u\) y \(h\). Así se cuentan bisagras distintas en H4.

\subsection{Conteo a dos escalas}\label{conteo-a-dos-escalas}

En un ciclo de longitud cuatro, los dos emparejamientos alternados deben usar anfitriones distintos. Al distinguir un emparejamiento hay \(r(r-1)\) elecciones, cada una con peso \(r^{-4}\) en el objetivo inicial. El multiplicador por ciclo es \((r-1)/r^3\). Hay \(\binom d2^2\) ciclos en \(K_{d,d}\) y \(3\binom d4\) en \(K_d\). Cuando \(r/d\to1/p\), las contribuciones divididas por \(d^2\) tienden a \(p^2/4\) y \(p^2/8\), respectivamente.

El núcleo original tiene 15 roles. Hay \(\binom{15}{2}=105\) pares de roles distintos y 15 tipos de clique en un solo rol. Con \(a=\lfloor j/2\rfloor\), los bloques tardíos y secundarios completos aportan a lo sumo \(108t^2(3/4)^a=o(t^2)\) bases de los tipos restringidos. El coeficiente principal para las bases del núcleo original es \[105\frac{p^2}{4}+15\frac{p^2}{8}=\frac{225p^2}{8}.\] Cada Pasch contado tiene dos anfitriones. Por el Lema 7, sus cuatro triángulos pertenecen a un único bloque. En la clique terminal \(w/n\to p^2/64\), así que su contribución dividida por \(n^2\) tiende a \((p^2/8)(p^2/64)^2=p^6/32768\).

Para cada exponente fijo \(q\), los bloques tempranos tienen lado al menos \(\sqrt t\), suficiente para que las \(n^q\) funciones de prueba formen una familia polinómica en el orden de cada bloque. T+ devuelve una sola descomposición por bloque para toda la familia y el conteo de Pasch. Al restringir a las bases de cada bloque, la unión es propia. La suma de errores en los bloques originales da (3), y la suma en todos los bloques da (2) para esa misma unión. La sucesión de reservas se elige antes de fijar las funciones de prueba y el borrado. La descomposición final puede depender de ambos.

\end{document}
